% !TEX root = ../main.tex

\section{Introducci\'on}
\label{sec:introduction}

Desde el nacimiento de la programaci\'on, la
interfaz a trav\'es de la cual los programas se materializan en c\'odigo
ha sido esencialmente la misma: un teclado y una pantalla. Llamamos a
este arreglo el \emph{paradigma de la m\'aquina de escribir}. Es una
herencia tan casual que rara vez se cuestiona, pero sus
consecuencias son profundas y de a ratos silenciosas. Condiciona qui\'enes pueden programar,
c\'omo aprenden a hacerlo y qu\'e formas de expresi\'on consideramos
correctas dentro de la computaci\'on.

Tres problemas convergen en este punto. Primero, la interfaz teclado-pantalla
discrimina contra usuarios cuyas capacidades f\'isicas, cognitivas o
sensoriales difieren del promedio asumido por sus dise\~nadores~\cite{ko23}.
Segundo, la programaci\'on impone una curva de aprendizaje pronunciada y
un umbral de entrada alto~\cite{bosse17gerosa}, lo que desalienta a una
parte importante de la poblaci\'on incluso antes de que llegue a escribir
su primera l\'inea de c\'odigo. Tercero, los lenguajes de programaci\'on 
contempor\'aneos y su docencia descansan sobre esta interfaz plana, trayendo consigo los dos problemas 
anteriores y sus repercusiones sobre quienes pretenden programar y c\'omo se acercan a la programaci\'on.
Si la intenci\'on es que la programaci\'on sea un medio universal y abierto para el pensamiento, entonces 
el paradigma de la m\'aquina de escribir es un obst\'aculo para pensar de formas distintas de lograrlo.
La soluci\'ion podr\'ia radicar en la intersecci\'on entre la computaci\'on y el arte~\cite{hongji16creative}.

La pregunta que orienta este trabajo es si es posible dise\~nar un
lenguaje de programaci\'on que (i) materialice los programas en un
soporte radicalmente distinto al texto plano, (ii) sirva y preserve las funcionalidades computacionales y las propiedades
formales que se esperan de un lenguaje  y (iii) brinde un ecosistema de herramientas creativas 
que permitan al usuario explorar e interactuar con la computaci\'on de forma libre y expresiva.
La hip\'otesis es entonces que la expresi\'on estetica puede servir 
como veh\'iculo legitimo para la computaci\'on. Ensamblar piezas tangibles puede ser, al tiempo,
un acto escult\'orico y un acto de computaci\'on.

\sculpt (\emph{Simple Cubic Language for Programming Tasks}) es la
implementaci\'on de esa hip\'otesis. En t\'erminos generales, \sculpt
es a la vez (a) un lenguaje de programaci\'on Turing-completo y su especificaci\'on formal, (b) una colecci\'on de piezas y sus caracteristicas de construcci\'on y
(c) \sculpter, un entorno y \emph{runtime} web escrito en Scala~3 y
Scala.js que permite escribir, depurar y ejecutar programas antes de
ensamblarlos f\'isicamente.

Este art\'iculo presenta el dise\~no fisico y l\'ogico de \sculpt, su sem\'antica, una
demostraci\'on formal condensada de su Turing-completitud y una
validaci\'on emp\'irica con varios grupos de participantes. La contribuci\'on principal
no es t\'ecnica en sentido estricto, sino conceptual: mostrar que la
\emph{tangibilidad}, la \emph{libertad est\'etica} y la
\emph{computaci\'on} pueden coexistir en un mismo lenguaje
sin que ninguna de las tres se sacrifique.

\endinput
